\begin{ThreePartTable}

    \renewcommand\TPTminimum{\textwidth}

    \setlength\LTleft{0pt}
    \setlength\LTright{0pt}
    \setlength\tabcolsep{3pt}

    \begin{TableNotes}
        \item \textbf{Nota.} Adaptado de: \textcite[p. 44]{AENOR:2016} [traducido del inglés].
    \end{TableNotes}

    \begin{longtable}{@{}P{0.2\textwidth}P{0.2\textwidth}P{0.6\textwidth} @{}}

        \caption{Datos de mantenimiento} \label{tab:maintenance_data} \\

        \toprule
        \textbf{Categoría de los datos} &
            \textbf{Datos a ser registrados} &
                \textbf{Descripción} \\
        \toprule
        \endfirsthead
        
        \multicolumn{3}{r}{\textit{(Continuación de la tabla~\ref{tab:maintenance_data})}} \\
        \toprule
        \textbf{Categoría de los datos} &
            \textbf{Datos a ser registrados} &
                \textbf{Descripción} \\
        \toprule
        \endhead

        \midrule[\heavyrulewidth]
        \multicolumn{3}{r}{\textit{Continúa en la siguiente página}}\\
        \endfoot

        \bottomrule
        \insertTableNotes\\
        \endlastfoot

        \multirow{3}{\linewidth}{Identificación} &
            Registro de mantenimiento &
                Identificación única del mantenimiento \\
            \cline{2-3} &
            Ubicación e identificación del equipo &
                p. ej. Etiqueta de identificación física \\
            \cline{2-3} &
            Registro de fallas &
                El correspondiente registro de identificación de falla \\

        \midrule

        \multirow{9}{\linewidth}{Datos de mantenimiento} & 
            Fecha del mantenimiento & 
                Fecha cuando la actividad de mantenimiento fue realizada o planificada (Fecha de inicio) \\
            \cline{2-3} &
            Ubicación e identificación del equipo &
                p. ej. Etiqueta de identificación física \\
            \cline{2-3} &
            Categoría del mantenimiento &
                Categoría principal (correctivo, preventivo) \\
            \cline{2-3} &
            Prioridad del mantenimiento &
                Prioridad alta, media o baja \\
            \cline{2-3} &
            Intervalo (mantenimiento planificado) &
                Calendario o intervalo operativo \\
            \cline{2-3} &
            Actividad de mantenimiento &
                Descripción de la actividad de mantenimiento \\
            \cline{2-3} &
            Impacto del mantenimiento en operaciones en planta &
                Cero, parcial o total \\
            \cline{2-3} &
            Sub-unidad mantenida &
                Nombre de la sub-unidad mantenida (puede omitirse para el mantenimiento preventivo) \\
            \cline{2-3} &
            Componentes o elementos mantenidos &
                Especificar el componente o elemento que fue mantenido \\
            \cline{2-3} &
            Localización de piezas de recambio &
                Disponibilidad de piezas de recambios (p. ej. local/distante, fabricante) \\

        \midrule

        \multirow{3}{\linewidth}{Recursos de mantenimiento} &
            Horas-hombre de mantenimiento por disciplina &
                Horas-hombre de mantenimiento por disciplina (mecánica, eléctrica, instrumentación, otros) \\
            \cline{2-3} &
            Horas-hombre de mantenimiento totales &
                Horas-hombre de mantenimiento totales \\
            \cline{2-3} &
            Recursos para equipos de mantenimiento &
                p. ej. recipiente para la intervención, grúa \\

        \midrule

        \multirow{3}{\linewidth}{Tiempos de mantenimiento} &
            Tiempo de mantenimiento activo &
                Tiempo de duración para que el mantenimiento activo sea realizado en el equipo \\
            \cline{2-3} & 
            Tiempo muerto &
                Tiempo de duración durante el cual un elemento está en inactividad \\
            \cline{2-3} &
            Problemas o retrasos de mantenimiento &
                Causas de tiempos de inactividad prolongados, p. ej. logísticas, climáticas, andamiaje, falta de repuestos, retraso del personal \\

        \midrule

        Observaciones &
            Información adicional &
                Proporcionar más detalles, si están disponibles, en la actividad de mantenimiento y recursos usados \\

    \end{longtable}

\end{ThreePartTable}